\documentclass[10pt, letterpaper]{article}

\usepackage[
    ignoreheadfoot,
    top=2 cm,
    bottom=2 cm,
    left=2 cm,
    right=2 cm,
    footskip=1.0 cm,
]{geometry}
\usepackage{titlesec}
\usepackage{tabularx}
\usepackage[T2A]{fontenc}
\usepackage[utf8]{inputenc}
\usepackage[russian,english]{babel}
\usepackage{array}
\usepackage[dvipsnames]{xcolor}
\definecolor{primaryColor}{RGB}{0, 0, 0}
\usepackage{enumitem}
\usepackage{fontawesome5}
\usepackage{amsmath}
\usepackage[
    pdftitle={Andrey Myuziev Resume},
    pdfauthor={Andrey Myuziev},
    pdfcreator={LaTeX with RenderCV},
    colorlinks=true,
    urlcolor=primaryColor
]{hyperref}
\usepackage[pscoord]{eso-pic}
\usepackage{calc}
\usepackage{bookmark}
\usepackage{lastpage}
\usepackage{changepage}
\usepackage{paracol}
\usepackage{ifthen}
\usepackage{needspace}
\usepackage{iftex}

\ifPDFTeX
    \input{glyphtounicode}
    \pdfgentounicode=1
    \usepackage{lmodern}
\fi

\raggedright
\AtBeginEnvironment{adjustwidth}{\partopsep0pt}
\pagestyle{empty}
\setcounter{secnumdepth}{0}
\setlength{\parindent}{0pt}
\setlength{\topskip}{0pt}
\setlength{\columnsep}{0.15cm}
\pagenumbering{gobble}

\titleformat{\section}{\needspace{4\baselineskip}\bfseries\large}{}{0pt}{}[\vspace{1pt}\titlerule]

\titlespacing{\section}{
    -1pt
}{
    0.3 cm
}{
    0.2 cm
}

\renewcommand\labelitemi{$\vcenter{\hbox{\small$\bullet$}}$}

\newenvironment{highlights}{
    \begin{itemize}[
        topsep=0.10 cm,
        parsep=0.10 cm,
        partopsep=0pt,
        itemsep=0pt,
        leftmargin=0 cm + 10pt
    ]
}{
    \end{itemize}
}

\newenvironment{onecolentry}{
    \begin{adjustwidth}{
        0 cm + 0.00001 cm
    }{
        0 cm + 0.00001 cm
    }
}{
    \end{adjustwidth}
}

\newenvironment{twocolentry}[2][]{
    \onecolentry
    \def\secondColumn{#2}
    \setcolumnwidth{\fill, 4.5 cm}
    \begin{paracol}{2}
}{
    \switchcolumn \raggedleft \secondColumn
    \end{paracol}
    \endonecolentry
}

\newenvironment{header}{
    \setlength{\topsep}{0pt}\par\kern\topsep\centering\linespread{1.5}
}{
    \par\kern\topsep
}

\let\hrefWithoutArrow\href

\begin{document}
    \newcommand{\AND}{\unskip
        \cleaders\copy\ANDbox\hskip\wd\ANDbox
        \ignorespaces
    }

    \newsavebox\ANDbox
    \sbox\ANDbox{$|$}

    \begin{header}
        \fontsize{25 pt}{25 pt}\selectfont Андрей Мюзиев

        \vspace{5 pt}

        \normalsize
        \mbox{Санкт-Петербург}%
        \kern 5.0 pt%
        \AND%
        \kern 5.0 pt%
        \mbox{\hrefWithoutArrow{mailto:andrej.myuziev@yandex.ru}{andrej.myuziev@yandex.ru}}%
        \kern 5.0 pt%
        \AND%
        \kern 5.0 pt%
        \mbox{\hrefWithoutArrow{tel:+79210144773}{+7 921 014 47 73}}%
        \kern 5.0 pt%
        \AND%
        \kern 5.0 pt%
        \mbox{\hrefWithoutArrow{https://github.com/Komperlan}{github.com/Komperlan}}%
    \end{header}

    \vspace{5 pt - 0.3 cm}

    \section{Профиль}

        \begin{onecolentry}
            Студент направления программной инженерии в ИТМО. Интересуюсь backend-разработкой и машинным обучением: проектировал микросервисные приложения на C\# и Java, собирал датасеты, обучал модели регрессии, кластеризации, аудио-классификации и генерации изображений.
        \end{onecolentry}

    \section{Образование}

        \begin{twocolentry}{
            сентябрь 2023 -- май 2027
        }
            \textbf{Университет ИТМО}, бакалавриат, программная инженерия
        \end{twocolentry}

    \section{Проекты}

        \begin{twocolentry}{
            \href{https://github.com/Komperlan/Labs/tree/main/ml/4-audio-Komperlan-main}{код}
        }
            \textbf{Классификация персонажей Dota 2 по голосовым репликам}
        \end{twocolentry}

        \vspace{0.10 cm}
        \begin{onecolentry}
            \begin{highlights}
                \item Собрал датасет из 1668 MP3-файлов для пяти персонажей, написал Selenium-парсер Dota 2 Fandom и подготовил mel-спектрограммы через librosa.
                \item Реализовал собственные LSTMCell и многослойный LSTM-классификатор на PyTorch; получил точность на тестовой выборке 0.9281 и визуализировал LSTM-эмбеддинги через t-SNE.
                \item Технологии: Python, PyTorch, librosa, scikit-learn, Selenium, BeautifulSoup, NumPy, Matplotlib.
            \end{highlights}
        \end{onecolentry}

        \vspace{0.2 cm}

        \begin{twocolentry}{
            \href{https://github.com/Komperlan/Labs/tree/main/ml/5-generative-Komperlan-main}{код}
        }
            \textbf{Генерация изображений с помощью вариационного автоэнкодера}
        \end{twocolentry}

        \vspace{0.10 cm}
        \begin{onecolentry}
            \begin{highlights}
                \item Обучил сверточный VAE для двух датасетов: 1554 иконки и 5813 anime-изображений, приводя изображения к 64x64 и работая с латентными пространствами размерности 512 и 256.
                \item Реализовал генерацию из случайного шума, восстановление изображений и GIF-интерполяции между объектами в латентном пространстве; ошибка на валидации для эксперимента с иконками снизилась до 0.0404.
                \item Технологии: Python, PyTorch, TorchVision, Torch DirectML, Pillow, imageio, NumPy, Matplotlib.
            \end{highlights}
        \end{onecolentry}

        \vspace{0.2 cm}

        \begin{twocolentry}{
            \href{https://github.com/Komperlan/Labs/tree/main/ml}{код}
        }
            \textbf{ML-пайплайн для анализа автомобильных объявлений Auto.ru}
        \end{twocolentry}

        \vspace{0.10 cm}
        \begin{onecolentry}
            \begin{highlights}
                \item Написал парсер Auto.ru, подготовил CSV с характеристиками автомобилей и построил пайплайн предобработки: очистка признаков, удаление выбросов, one-hot-кодирование и стандартизация.
                \item Для предсказания цены сравнил базовую модель, HistGradientBoostingRegressor и собственный градиентный бустинг; лучшая sklearn-модель получила RMSE на тесте 1 002 487 и R2 0.704.
                \item Для кластеризации сравнил базовое разбиение по цене, sklearn K-Means и собственный K-Means; собственная реализация достигла Silhouette Score 0.2148.
                \item Технологии: Python, pandas, NumPy, scikit-learn, BeautifulSoup, requests, Matplotlib.
            \end{highlights}
        \end{onecolentry}

        \vspace{0.2 cm}

        \begin{twocolentry}{
            \href{https://github.com/Komperlan/Labs/blob/main/c\%23\%20microservice/lab-3/lab-3.md}{код}
        }
            \textbf{C\# веб-приложение с гексагональной архитектурой}
        \end{twocolentry}

        \vspace{0.10 cm}
        \begin{onecolentry}
            \begin{highlights}
                \item Разработал микросервисную систему с изолированной бизнес-логикой, gRPC API, HTTP-шлюзом, хранением данных в PostgreSQL, аудитом операций и транзакционным сервисным слоем.
                \item Реализовал репозитории с пагинацией и фильтрацией, middleware для обработки ошибок, серверные interceptors и полиморфные модели через Protocol Buffers oneof.
                \item Технологии: ASP.NET Core, C\#, gRPC, PostgreSQL, Npgsql, FluentMigrator, Swagger.
            \end{highlights}
        \end{onecolentry}

        \vspace{0.2 cm}

        \begin{twocolentry}{
            \href{https://github.com/Komperlan/Labs/blob/main/Java/Lab4/lab-4.md}{код}
        }
            \textbf{Java веб-приложение на Spring Boot}
        \end{twocolentry}

        \vspace{0.10 cm}
        \begin{onecolentry}
            \begin{highlights}
                \item Разработал Spring Boot приложение с REST API поверх существующей PostgreSQL базы через Spring Data JPA и Hibernate.
                \item Добавил аутентификацию и правила авторизации: владельцы и администраторы могут изменять данные, все авторизованные пользователи могут читать.
                \item Технологии: Java, Spring Boot, Spring Data JPA, Hibernate, PostgreSQL, SQL, Gradle.
            \end{highlights}
        \end{onecolentry}

    \section{Технологии}

        \begin{onecolentry}
            \textbf{Языки:} Python, C++, Java, C\#, SQL, Protobuf
        \end{onecolentry}

        \vspace{0.2 cm}

        \begin{onecolentry}
            \textbf{ML и данные:} PyTorch, scikit-learn, pandas, NumPy, librosa, Matplotlib, BeautifulSoup, Selenium
        \end{onecolentry}

        \vspace{0.2 cm}

        \begin{onecolentry}
            \textbf{Backend:} Spring, PostgreSQL, Linux, Docker, .NET, Gradle, gRPC, HTTP, Hibernate
        \end{onecolentry}

\end{document}
